\chapter{Apnea}

\section*{Definition}
Apnea is a pause in breathing. In adult sleep testing, an apnea is generally scored when airflow stops for at least 10 seconds; in infants, a pause of 20 seconds or a shorter pause with bradycardia, cyanosis, pallor, or oxygen desaturation is clinically significant. Central apnea occurs without respiratory effort, obstructive apnea occurs despite effort against a blocked airway, and mixed apnea has both patterns.

\section*{When to See a Doctor}

\begin{itemize}
 \item Witnessed apnea episodes with cyanosis, loss of consciousness, or seizure-like activity
 \item Excessive daytime sleepiness, loud snoring, or witnessed breathing pauses
 \item Gasping or choking sensations during sleep reported by bed partner
 \item Any infant or child with a breathing pause, color change, limpness, or need for stimulation or resuscitation
 \item Morning headaches, cognitive impairment, or mood disturbances suggestive of chronic sleep disruption
\end{itemize}

\section*{How It Develops}
Obstructive apnea occurs when the upper airway narrows or collapses during sleep despite continued effort to breathe. Central apnea occurs when the brain temporarily reduces or stops respiratory drive. Some episodes combine both mechanisms. The cause and significance differ between adults, children, and newborns.

\section*{Causes}
\begin{itemize}
 \item \textbf{Obstructive sleep apnea:} Obesity, retrognathia, macroglossia, tonsillar hypertrophy, nasal obstruction, neuromuscular disorders
 \item \textbf{Central sleep apnea:} Heart failure (Cheyne-Stokes respiration), stroke, opioid use, high-altitude periodic breathing, brainstem lesions
 \item \textbf{Neonatal apnea:} Prematurity and immature respiratory control, infection, intracranial disease, metabolic disorders, or medication exposure; gastroesophageal reflux is not usually considered a cause of apnea
 \item \textbf{Transient:} Breath-holding spells (children), voluntary apnea, post-hyperventilation
\end{itemize}

\section*{Related Symptoms}
\begin{itemize}
 \item Snoring, daytime hypersomnolence, morning headache, nocturia, dry mouth upon waking, irritability, witnessed gasping
\end{itemize}
\textbf{Important:} Untreated obstructive sleep apnea is associated with sleepiness, impaired attention, hypertension, and cardiovascular risk. A diagnosis requires appropriate sleep testing rather than symptoms alone.

\section*{Medical Evaluation}
\begin{itemize}
 \item Polysomnography or, for uncomplicated adults at increased risk of moderate-to-severe obstructive sleep apnea, a clinician-ordered home sleep apnea test. Children and people with significant cardiopulmonary, neuromuscular, opioid-related, or other complex disease generally need in-lab assessment.
 \item Apnea-hypopnea index (AHI) is interpreted with age, symptoms, and comorbidities. Adult categories of 5--less than 15, 15--less than 30, and 30 or more events per hour are not pediatric thresholds.
 \item Oxygen desaturation index and nocturnal oxygen nadir assessment
 \item Electrocardiogram, echocardiogram if suspicion of heart failure as cause of central apnea
 \item Brain imaging (MRI) if central apnea with neurologic signs or suspected brainstem lesion
\end{itemize}

\section*{Treatment Options}
\begin{itemize}
 \item Positive airway pressure, usually CPAP or auto-adjusting PAP, is a common effective treatment for adult obstructive sleep apnea when prescribed after diagnosis
 \item Mandibular advancement device for mild-to-moderate OSA or CPAP intolerance
 \item Central sleep apnea is treated according to its cause, which may include treating heart failure, changing opioid therapy under supervision, CPAP, other PAP modes, oxygen, or specialist therapies. Adaptive servo-ventilation requires specialist selection and close monitoring, particularly in heart failure with reduced ejection fraction.
 \item Upper airway surgery (uvulopalatopharyngoplasty, tonsillectomy) for surgically correctable obstruction
 \item Weight-management treatment, including bariatric surgery for selected people with obesity, may improve but does not automatically cure obstructive sleep apnea; follow-up testing is needed.
\end{itemize}

\section*{Self-Care and Home Management}
\begin{itemize}
 \item Weight loss for overweight/obese individuals with sleep apnea (even 10\% reduction can improve symptoms)
 \item Positional therapy: sleep on side rather than supine to reduce airway collapse
 \item Avoid alcohol and sedatives near bedtime unless prescribed; never stop prescribed sedatives or opioids abruptly without medical advice
 \item Positional therapy or modest head elevation may help some people, but should not replace prescribed PAP treatment
 \item Nasal saline irrigation and allergy management if nasal congestion contributes to obstruction
\end{itemize}

\section*{References}
\begin{thebibliography}{99}
\bibitem{apnea_general:ref1} Jordan AS, McSharry DG, et al. ``Adult obstructive sleep apnoea.'' \textit{Lancet}. 2014;383(9918):736-747.
\bibitem{apnea_general:ref2} Epstein LJ, Kristo D, et al. ``Clinical guideline for the evaluation, management and long-term care of obstructive sleep apnea in adults.'' \textit{J Clin Sleep Med}. 2009;5(3):263-276.
\bibitem{apnea_general:ref3} Eckert DJ, Malhotra A, et al. ``Pathophysiology of adult obstructive sleep apnea.'' \textit{Proc Am Thorac Soc}. 2008;5(2):144-153.
\bibitem{apnea_general:ref4} Kryger MH, Roth T, et al. \textit{Principles and Practice of Sleep Medicine}, 7th ed. Elsevier, 2022.
\bibitem{apnea_general:ref5} UpToDate. ``Clinical presentation and diagnosis of obstructive sleep apnea in adults.'' Wolters Kluwer, 2023.
\bibitem{apnea_general:ref6} American Academy of Sleep Medicine. Clinical practice guideline for diagnostic testing for adult obstructive sleep apnea. 2017.
\bibitem{apnea_general:ref7} American Academy of Sleep Medicine. Treatment of central sleep apnea in adults: clinical practice guideline. \textit{J Clin Sleep Med}. 2025;21(12):2213--2236.
\end{thebibliography}
